% !TeX root = ../main.tex
% ========== Sperrvermerk ==========
% Text ggf. an die aktuelle Vorgabe von Firma/Hochschule anpassen.
\section*{Sperrvermerk}
\addcontentsline{toc}{section}{Sperrvermerk}
\thispagestyle{plain}

Die vorliegende \ArbeitTyp{} wurde im Rahmen eines laufenden Projektes bei der Firma \Firma{} durchgeführt. Zur Wahrung der Vertraulichkeit verpflichte ich mich, von Beginn an bis mindestens fünf Jahre nach Abgabe der \ArbeitTyp{}, Veröffentlichungen jeglicher Art der \ArbeitTyp{} oder ihres wesentlichen Inhalts nicht vorzunehmen, sofern keine ausdrückliche schriftliche Freigabe durch \FirmaKurz{} erfolgt. Aus diesem Grund darf während des genannten Zeitraums auch kein Exemplar der \ArbeitTyp{} zur Einsicht oder Ausleihe in der Bibliothek der \Hochschule{} aufgenommen werden, es sei denn, eine entsprechende Freigabe durch die Firma \FirmaKurz{} liegt vor.

\vspace{1em}
\noindent
Ich erkläre mich damit einverstanden, dass das Thema meiner \ArbeitTyp{} unter Nennung meines Namens in einem Verzeichnis der Hochschule aufgeführt wird.

\vspace{1em}

\begin{tabular}{p{6cm} p{6cm}}
	Ort, Datum: & Unterschrift: \\[1.5em]
	\hrulefill & \hrulefill
\end{tabular}

\vspace{2em}
\noindent
\textbf{Von der Hochschule zu unterschreiben}

\vspace{1em}
\noindent
Aus den oben genannten Gründen verpflichtet sich die Hochschule, Veröffentlichungen der \ArbeitTyp{} oder ihres wesentlichen Inhalts nicht vorzunehmen und die Arbeit nicht zur Einsicht oder Ausleihe in die Bibliothek aufzunehmen, sofern keine ausdrückliche schriftliche Freigabe durch \FirmaKurz{} vorliegt.

\vspace{1em}

\begin{tabular}{p{6cm} p{6cm}}
	Ort, Datum: & Unterschrift: \\[1.5em]
	\hrulefill & \hrulefill
\end{tabular}

\vspace{2em}
\noindent
\textbf{Vom Betreuer der \FirmaKurz{} zu unterschreiben}

\vspace{1em}
\noindent
Der oben genannte Sperrvermerk wurde zur Kenntnis genommen.

\vspace{1em}

\begin{tabular}{p{6cm} p{6cm}}
	Ort, Datum: & Unterschrift: \\[1.5em]
	\hrulefill & \hrulefill
\end{tabular}
